% !PRINTED 659
\weberpage{659}

%<<CONT
\weq{1}{\eta = \frac{\mu}{\nu},}

und $\mu$ heißt der Zähler, $\nu$ der Nenner von $\eta$. Diese beiden
Funktionen sind durch $\eta$ nicht vollständig bestimmt, wohl aber
sind die Funktionale bestimmt, die übrig bleiben, wenn alle
gemeinschaftlichen Funktionalfaktoren im Zähler und Nenner
herausgehoben werden. Denn ist

\weq{2}{\frac{\mu}{\nu} = \frac{\mu'}{\nu'}; \quad \mu\nu' = \nu\mu',}

so muß jedes Primfunktional, das in $\nu$, aber nicht in $\mu$
aufgeht, in $\nu'$ aufgehen. Die so aus $\eta$ bestimmten Funktionale
$\bar\alpha$ und $\bar\beta$ nennen wir Zählerfunktional und
Nennerfunktional von $\eta$. Um $\eta$ durch Funktionen darzustellen,
nehme man zunächst eine Funktion $\nu$, teilbar durch $\bar\beta$,
aber sonst beliebig und setze $\nu = \bar\beta\bar\gamma$; darin kann
$\bar\gamma$ relativ prim zu einem beliebigen Funktional angenommen
werden. Dann ist $\eta\nu$ eine durch $\bar\gamma$ teilbare ganze
Funktion, und man setze

\[ \eta\nu = \mu = \bar\alpha\bar\gamma. \]

Das Funktional $\bar\alpha$ kann aber, wie wir später noch
nachweisen werden, bei gegebenem $\bar\beta$ nicht mehr beliebig
sein.

Die Funktionale $\bar\alpha, \bar\beta$ werden wir auch kurz \gesperrt{den
Zähler und den Nenner} von $\eta$ nennen.

Ist $\pi$ ein Primfunktional, so gibt es \gesperrt{Zahlen} (Konstanten) $a, b,
a', b'$, die den Kongruenzen

\[ \mu \equiv a,\ \ \nu \equiv b,\ \ \mu' \equiv a',\ \ \nu' \equiv b' \quad (\mathrm{mod}\ \pi) \]

genügen, und aus (2) folgt

\[ ab' = ba', \quad \frac{a}{b} = \frac{a'}{b'} = c, \]

vorausgesetzt, daß $\nu$ und $\nu'$ nicht durch $\pi$ teilbar sind.
Es ist dann $\eta - c$ eine Funktion in $\Omega$, in der der Zähler,
aber nicht der Nenner durch $\pi$ teilbar ist. Ist $\varrho$ eine
ganze oder gebrochene Funktion, in der der Zähler durch $\pi$, aber
nicht durch $\pi^2$, der Nenner auch nicht durch $\pi$ teilbar ist,
so ist $(\eta - c) : \varrho$ eine Funktion, deren Nenner nicht durch
$\pi$ teilbar ist und deren Zähler den Faktor $\pi$ einmal weniger
enthält als $(\eta - c)$. Wir setzen:

\[ \eta - c = \varrho\eta_1. \]

Dieselbe Betrachtung wenden wir auf $\eta_1$ an und bekommen so eine
beliebig fortzusetzende Reihe von Gleichungen:
%CONT>>
